\documentclass[11pt]{article}

\usepackage[margin=1in]{geometry}
\usepackage[T1]{fontenc}
\usepackage[utf8]{inputenc}
\usepackage{graphicx}
\usepackage{booktabs}
\usepackage{tabularx}
\usepackage{array}
\usepackage{hyperref}
\usepackage{xcolor}
\usepackage{enumitem}
\usepackage{float}

\hypersetup{
  colorlinks=true,
  linkcolor=blue,
  urlcolor=blue
}

\newcommand{\ok}{\textcolor{green!50!black}{Implemented}}
\newcommand{\partialstatus}{\textcolor{orange!85!black}{Partial}}
\newcommand{\missingstatus}{\textcolor{red!70!black}{Missing}}

\title{GPU Request Draft:\\PikoGPT Pre-Training Readiness}
\author{FunkyAI}
\date{April 7, 2026}

\begin{document}
\maketitle

\section*{Proposed Model Architectures}

\subsection*{Main Configuration: \texttt{train\_large.toml}}
Our main proposed model for the GPU run is \texttt{train\_large.toml}. This is the strongest configuration currently defined in the repository while remaining within the 40M-parameter budget.

\begin{itemize}[leftmargin=1.5em]
  \item Architecture: decoder-only transformer
  \item Embedding dimension: 384
  \item Transformer layers: 10
  \item Attention heads: 6
  \item Context length: 512
  \item Dropout: 0.1
  \item Total parameters: 37.24M
\end{itemize}

This model is our preferred choice because it offers the best balance of model capacity, context length, and training practicality for an 8-GPU run.

\subsection*{Plan B: \texttt{train\_deep.toml}}
\texttt{train\_deep.toml} is our first fallback option. It keeps the 512-token context length but shifts capacity toward depth instead of width.

\begin{itemize}[leftmargin=1.5em]
  \item Architecture: decoder-only transformer
  \item Embedding dimension: 320
  \item Transformer layers: 14
  \item Attention heads: 8
  \item Context length: 512
  \item Dropout: 0.1
  \item Total parameters: 33.51M
\end{itemize}

We would use this configuration if a deeper model appears more stable or more sample-efficient than the main configuration while keeping the same context length.

\subsection*{Plan B: \texttt{train\_fullcontext.toml}}
\texttt{train\_fullcontext.toml} is our second fallback option. It prioritizes the full 1024-token context window while remaining comfortably below the 40M-parameter limit.

\begin{itemize}[leftmargin=1.5em]
  \item Architecture: decoder-only transformer
  \item Embedding dimension: 352
  \item Transformer layers: 10
  \item Attention heads: 8
  \item Context length: 1024
  \item Dropout: 0.1
  \item Total parameters: 32.97M
\end{itemize}

We would use this configuration if longer-range context becomes more valuable than maximizing parameter count.

\section*{Full Training Command}
The repository already contains a full pipeline script that performs preprocessing, distributed training, and final inference. For the main GPU run we plan to use the following command:

\begin{verbatim}
SOURCE_DATASET_PATH=data/raw/openwebtext_local \
TEST_DATA_PATH=data/raw/NLP26_OWT_eval/test \
NPROC_PER_NODE=8 \
RUN_NAME=pikogpt_large_gpu_run \
CONFIG=configs/train_large.toml \
scripts/full_training_run.sh
\end{verbatim}

\section*{Checklist Status}
\begin{table}[H]
\centering
\small
\begin{tabularx}{\textwidth}{>{\raggedright\arraybackslash}p{0.40\textwidth} >{\raggedright\arraybackslash}p{0.18\textwidth} X}
\toprule
Item & Status & Evidence / comment \\
\midrule
Technical requirements & \ok & Decoder-only model, GPT-2 tokenizer, model configs below 40M parameters, config-driven training pipeline. \\
Checkpointing and logging & \ok & Periodic step checkpoints, resumable training, debug logs, train/eval JSONL metrics. \\
Full training pipeline and script & \ok & \texttt{main.py} stages plus \texttt{scripts/full\_training\_run.sh}. \\
Plot and evaluation pipeline & \partialstatus & Validation-loss evaluation exists during training; plotting is now supported via \texttt{scripts/plot\_training\_metrics.py}; no benchmark suite stage yet. \\
Distributed training on 8 GPUs & \partialstatus & DDP support is implemented and verified locally with 2 processes; 8-GPU execution is not directly testable on this machine. \\
Decisions and results outlined & \partialstatus & Hyperparameter tuning results exist; architecture variants exist; architecture-comparison run results are still incomplete. \\
\bottomrule
\end{tabularx}
\end{table}

\section*{Smoke Results Evaluation}

\subsection*{Smoke Run Overview}
We ran smoke validations to confirm that the full infrastructure works before requesting large-scale GPU time. These smoke runs are not meant to demonstrate final model quality; they are infrastructure checks for preprocessing, training, checkpointing, evaluation logging, resume support, and distributed execution.


\subsection*{Smoke Run Graphs}
\begin{figure}[H]
\centering
\includegraphics[width=0.32\textwidth]{gradient_norm_curve.png}
\includegraphics[width=0.32\textwidth]{learning_rate_curve.png}
\includegraphics[width=0.32\textwidth]{loss_curve.png}
\caption{Distributed smoke run diagnostics for test run: loss curve, learning-rate, and gradient norm. These plots were generated from the JSONL training logs.}
\end{figure}

\section*{Optuna Results Evaluation}

\end{document}
